\documentclass[10pt,letterpaper]{article}

\usepackage[T1]{fontenc}
\usepackage[utf8]{inputenc}
\usepackage{lmodern}
\usepackage[margin=0.90in]{geometry}
\usepackage{microtype}
\usepackage{enumitem}
\usepackage{xurl}
\usepackage[hidelinks]{hyperref}
\usepackage{titlesec}

\hypersetup{
  pdftitle={Qayd: Yet Another Constraint Solver with Lazy Clause Generation, Written in Rust},
  pdfauthor={Sohaib Afifi (Univ. Artois, UR 3926, LGI2A)},
  pdfsubject={Solver description for the 2026 XCSP3 Competition},
  pdfkeywords={Qayd, XCSP3, constraint programming, lazy clause generation}
}

\pagestyle{empty}
\setlength{\parindent}{0pt}
\setlength{\parskip}{0.34em}
\setlist[itemize]{leftmargin=1.25em,itemsep=0.08em,topsep=0.15em,parsep=0pt}
\titleformat{\section}{\large\bfseries}{\thesection}{0.55em}{}
\titlespacing*{\section}{0pt}{0.65em}{0.20em}

\title{\vspace{-1.8em}\textbf{Qayd: Yet Another Constraint Solver\\
with Lazy Clause Generation, Written in Rust}\vspace{-0.4em}}
\author{Sohaib Afifi\\[-0.15em]
\small Univ. Artois, UR 3926, Laboratoire de G\'enie Informatique et\\[-0.15em]
\small d'Automatique de l'Artois (LGI2A), F-62400 B\'ethune, France\\[-0.15em]
\small \href{mailto:sohaib.lafifi@univ-artois.fr}{sohaib.lafifi@univ-artois.fr}}
\date{\small Version 0.8.0, June 2026\vspace{-1.0em}}

\begin{document}
\maketitle
\thispagestyle{empty}
\vspace{-1.3em}

\begin{abstract}
\noindent
Qayd is an open-source constraint-programming solver written in Rust for
constraint satisfaction problems and mono-objective constraint optimization
problems over finite-domain integer variables. The v0.8.0 release targets the
XCSP$^3$-core catalogue and combines event-driven finite-domain propagation
with Lazy Clause Generation (LCG). Complete modes use conflict-driven search
and branch-and-bound, while parallel and incumbent-oriented modes add
portfolio search, selective information sharing, Large Neighborhood Search,
and incremental local search. This note describes the solver artifact submitted
to the 2026 XCSP$^3$ Competition.
\end{abstract}

\section{Frontend and propagation kernel}

The XCSP$^3$ frontend is built on a revision-pinned Rust parser and posts each
accepted callback directly into the solver. It reads plain XML as well as
\texttt{.lzma} and \texttt{.xz} files, handles CSP and mono-objective COP
instances, and emits the competition status, objective, and instantiation
records. The supported catalogue is a broad subset of XCSP$^3$-core
\cite{xcsp3,xcsp3core}. It includes intensional and extensional constraints,
\texttt{regular}, \texttt{mdd}, comparison and counting globals, extrema and
\texttt{element}, channeling, packing, scheduling, \texttt{circuit},
\texttt{instantiation}, and parser-expanded \texttt{slide}. Unsupported shapes
are rejected explicitly. In particular, product and lexicographic objectives
are outside the v0.8.0 frontend.

Domains are selected from three trailed representations: a dense sparse-set for
compact ranges, an explicit-support sparse-set for sparse domains, and
reversible bounds with trailed interior holes for very wide ranges. Propagators
subscribe at fix, bound-change, or domain-change granularity. A FIFO queue and
duplicate suppression drive propagation to a fixpoint, and root propagation
removes fixed non-objective variables before search.

Qayd provides dedicated filtering for several important globals. Positive
tables use Compact-Table-style reversible bitsets and residues
\cite{compacttable}. \texttt{regular} and \texttt{mdd} use layered
forward-backward reachability, following the automaton view of
\cite{pesant}. \texttt{allDifferent} enforces R\'egin's matching-based domain
consistency \cite{regin}. Fixed-resource \texttt{cumulative} combines
time-tabling, overload checks, and edge-finding. Other accepted forms are
implemented by dedicated propagators or sound decompositions.

\section{Learning and complete search}

The exact search layer follows LCG \cite{lcg}. Integer-domain facts are exposed
as equality atoms $[x=v]$ and order atoms $[x\geq k]$. Atoms for sparse domains
are allocated only for supported values, while atoms for wide ranges are
created on demand. Propagators attach explicit premises to domain changes when
available; otherwise the engine constructs a sound explanation from the
pre-propagation scope. Clausal propagation uses two watched literals and is
interleaved with finite-domain propagation.

Conflicts are analyzed at the first unique implication point. Qayd learns an
asserting clause, backjumps non-chronologically, and maintains variable activity
and Literal Block Distance (LBD). Search combines dom/wdeg
\cite{domwdeg} with VSIDS activity by minimizing a domain-size to combined-score
ratio. Phase saving, seed-dependent endpoint decisions, Luby restarts, an
LBD-degradation restart trigger \cite{lbd}, learned-clause reduction, bounded
root probing, and short CP-aware clause vivification complete the default
configuration. A lone sequential search periodically inverts its default
polarity for one restart segment; portfolio workers instead obtain diversity
from their seeds and restart schedules.

A single-worker CSP run uses CDCL to find one solution or prove
unsatisfiability. COP uses CDCL branch-and-bound and posts a strictly improving
bound after each incumbent. The objective can remain a variable, a symbolic
linear form, or an expression, which avoids materializing wide objective
ranges solely for search.

\section{Parallel and incumbent-oriented modes}

Multiworker CSP solving is heterogeneous: one worker performs chronological
depth-first search, while the remaining workers use diversified CDCL. CDCL
workers exchange only learned clauses of length at most eight and LBD at most
four. The first solution or complete unsatisfiability proof terminates the
portfolio.

Parallel COP workers share the incumbent value atomically. When the objective
is materialized as a variable, proof workers may also exchange short low-LBD
clauses. Additional roles include disjoint proof-space splitting, optimistic
objective probes, adaptive incumbent-driven Large Neighborhood Search
\cite{lns}, and a local-search incumbent generator. Heuristic incumbents used by
the proof portfolio are checked by fixing them in a cloned finite-domain model
and propagating before they update the shared bound.

The incumbent-only modes never report an optimality proof. With one worker,
\texttt{-{}-fast-cop} uses the exact engine with objective-biased value choices;
with several workers it dispatches local search. \texttt{-{}-turbo} always uses
local search and enables Guided Local Search weights together with a bounded
min-conflicts candidate set for large domains. The local-search engine maintains
incremental violation scores, uses constructive starts for selected table and
element patterns, repairs Boolean exact-cover rows preferentially, and applies
randomized restarts.

\section{Competition interface and availability}

The common invocation is:
\begin{center}
\small\ttfamily
qayd -t SECONDS -{}-seed SEED -p THREADS [mode] instance.xml[.lzma|.xz]
\end{center}
where \texttt{mode} can select \texttt{-{}-fast-cop}, \texttt{-{}-turbo}, or
optional COP roles such as \texttt{-{}-split}, \texttt{-{}-probe}, and
\texttt{-{}-lns}. If explicit values are absent, Qayd reads \texttt{RANDOMSEED}
and \texttt{NBCORE}. The release is identified by Git tag \texttt{v0.8.0} and
commit \texttt{e72a9e10c8d8c01ce37b5ae2ec379f78f322ff42}. Source code and platform
binaries are available at \url{https://github.com/sohaibafifi/qayd/releases/tag/v0.8.0}
under the MIT License.

\vspace{-0.25em}
\begin{thebibliography}{9}
\small
\setlength{\itemsep}{0.05em}
\setlength{\parskip}{0pt}

\bibitem{xcsp3}
F. Boussemart, C. Lecoutre, G. Audemard, and C. Piette.
XCSP3: An integrated format for benchmarking combinatorial constrained
problems. \emph{arXiv:1611.03398}, 2016.

\bibitem{xcsp3core}
F. Boussemart, C. Lecoutre, G. Audemard, and C. Piette.
XCSP3-core: A format for representing constraint satisfaction and optimization
problems. \emph{arXiv:2009.00514}, 2020.

\bibitem{compacttable}
J. Demeulenaere, R. Hartert, C. Lecoutre, G. Perez, L. Perron,
J.-C. R\'egin, and P. Schaus. Compact-Table: Efficiently filtering table
constraints with reversible sparse bit-sets. In \emph{CP 2016}, pages 207-223.

\bibitem{pesant}
G. Pesant. A regular language membership constraint for finite sequences of
variables. In \emph{CP 2004}, pages 482-495.

\bibitem{regin}
J.-C. R\'egin. A filtering algorithm for constraints of difference in CSPs.
In \emph{AAAI 1994}, pages 362-367.

\bibitem{lcg}
T. Feydy and P. J. Stuckey. Lazy clause generation reengineered.
In \emph{CP 2009}, pages 352-366.

\bibitem{domwdeg}
F. Boussemart, F. Hemery, C. Lecoutre, and L. Sa\"is.
Boosting systematic search by weighting constraints.
In \emph{ECAI 2004}, pages 146-150.

\bibitem{lbd}
G. Audemard and L. Simon. Predicting learnt clauses quality in modern SAT
solvers. In \emph{IJCAI 2009}, pages 399-404.

\bibitem{lns}
P. Shaw. Using constraint programming and local search methods to solve vehicle
routing problems. In \emph{CP 1998}, pages 417-431.

\end{thebibliography}

\end{document}
